\documentclass[french]{article}

\usepackage{babel}
\usepackage{comment}
\usepackage{psfig}
\newenvironment{solution}{\nopagebreak\par\vspace*{2mm}\underline{Solution~:}\em\par}{\par\vspace*{2mm}}
\usepackage{times}
\usepackage[latin1]{inputenc}


\pagestyle{headings}
\markright{28/6/1995 \hfill Conservatoire National des Arts et
M\'etiers -- Paris \hfill\ } 

\baselineskip 6pt
\parskip 6pt
\onecolumn
\textwidth 16cm
\textheight 24cm
\hoffset=-2cm
\voffset=-2.5cm
\setlength{\parindent}{0.0cm}
\unitlength 1cm
%\input{psfig}
\newtheorem{defin}{D\'efinition}[section]
\newtheorem{exampl}{Exemple}[section]
\def\picwidth{16cm}

%\excludecomment{solution}

\begin{document}
\begin{center}
\Large\bf Processus d'Informatisation et Bases de Donn\'ees A \\ 
    Examen Processus d'Informatisation\\
\large\bf Samedi 16 septembre 2000
\end{center}



On s'int�resse � l'informatisation d'un zoo. Ce zoo g�re
un parc d'animaux, ces animaux occupent des emplacements
et les employ�s du zoo ont en charge l'entretien 
des emplacements et du zoo. On vous demande 
de d\'ecrire \`a l'aide des outils vus en cours et en TD 
les principaux composants de cette application.


\begin{enumerate}
   \item Le zoo est constitu� d'un ensemble d'emplacements o�
	sont enferm�s les animaux. Chaque emplacement est pris
	en charge par un ou plusieurs employ�s, et un employ�
	peut s'occuper de plusieurs emplacements.

	\begin{figure}[htb]
	\centerline{\psfig{figure=../PI/exPI97-sept1.eps}}
	\caption{Le zoo}
	\label{exPI97-sept1}
	\end{figure}

       Compl\'etez le mod\`ele de la figure~\ref{exPI97-sept1},
	en ajoutant les cardinalit\'es. ({\bf 2 pts}).

   \item Chaque employ� s'occupe d'un emplacement � un
	horaire donn�. O� placer cette information {\it horaire}
	 dans le sch�ma pr�c�dent\,? {\bf (2 pts)}

   \item On veut maintenant r\'ef\'erencer le type d'entretient
	  effectu� par un employ� sur un emplacement.


	\begin{figure}[htb]
	\centerline{\psfig{figure=../PI/exPI97-sept2.eps}}	
	\caption{Entretien des emplacements}
	\label{exPI97-sept2}
	\end{figure}

        Voici (figure~\ref{exPI97-sept2}) une proposition
	pour introduire le type d'entretien. Vous para\^{\i}t-elle
	correcte\,? Si oui, quelles sont les cardinalit\'es.
	({\bf 2 pts}).

     \item Transformez l'association ``Entretient'' dans la 
	figure~\ref{exPI97-sept2} en entit\'e. Donnez le
	nouveau mod\`ele, sans oublier les cardinalit\'es. 
	({\bf 2 pts}). 

     \item Introduisez maintenant les informations suivantes
	   dans le mod�le de donn�es\,: {\bf (4 pts)} 
	 \begin{itemize}
	   \item Un des employ�s
	   dirige le zoo (un seul directeur, qui dirige un seul zoo).
	    \item Sur chaque emplacement, il y a un ou plusieurs
	      animaux (un animal est sur un seul emplacement\,!).
            \item Chaque animal appartient � une esp�ce (les ours, 
		     les tigres, ...)
         \end{itemize}
    
     \item Maintenant on s'int�resse aux relations de descendance
	    entre animaux\,: on souhaite savoir qu'un animal
	    est l'enfant d'un couple d'animaux. Comment
	    repr�senter cette information\,? {\bf (2 pts)}

   \item Donner le sch�ma relationnel (nom des tables,
     des attributs, cl�s primaires soulign�es) obtenu en rassemblant
	   les figures~\ref{exPI97-sept1} et \ref{exPI97-sept2}. 
	   Expliquer  la d�marche suivie.

\end{enumerate}
 


    
\end{document}
